% ============================================================================
% Draft section for Paper/paper.tex  (Agent S, discussion2, round 2)
% Suggested placement: after \section{The Symmetric Combination $m(G,H)$}
% (label sec:metric) and before \section{Computational Results}.
% Notation follows the paper: |T,G| is the non-homomorphism factor,
% \alpha is the INDEPENDENCE number (do not use it for |T,G|),
% \mu(f), \gamma(G,H) as in Definition~\ref{def:gamma}.
% New bibitems are at the end of this file; paste them into thebibliography.
% Items marked  % [check]  were cited from memory and not verified online.
% ============================================================================

\section{The Profile of a Graph}
\label{sec:profile}

The quantity $|T,G|$ can also be read with $G$ as the \emph{target}.
For a fixed $G$, it then records how badly each test graph $T$ fails to
map into $G$.  Normalising by $|E(T)|$ makes graphs of different sizes
comparable.

Throughout, $\operatorname{mc}_G(T)=|E(T)|-|T,G|$ denotes the largest number
of edges of $T$ preserved by a vertex map $V(T)\to V(G)$.  This is the
optimum of the \textsc{Max $G$-Col} problem~\cite{fjt2009,efjt2015}, also
called maximum graph homomorphism~\cite{lrs2006}.

\begin{definition}[Profile]
\label{def:profile}
Let $\mathfrak F$ be a countable family of graphs, each with at least one
edge (the \emph{test graphs}).  For a graph $G$ with at least one edge, the
\emph{profile} of $G$ is
\[
  \varphi(G)=\bigl(\varphi_T(G)\bigr)_{T\in\mathfrak F},\qquad
  \varphi_T(G)=\frac{|T,G|}{|E(T)|}\in[0,1].
\]
Given weights $w\colon\mathfrak F\to(0,\infty)$ with
$\sum_T w_T<\infty$, the \emph{profile distances} are
\[
  E_w(G,H)=\Bigl(\sum_{T\in\mathfrak F}w_T\bigl(\varphi_T(G)-\varphi_T(H)\bigr)^2\Bigr)^{1/2},
  \qquad
  E_\infty(G,H)=\sup_{T\in\mathfrak F}\bigl|\varphi_T(G)-\varphi_T(H)\bigr|.
\]
\end{definition}

The profile is the target-side, max-plus analogue of Lov\'asz's vector of
homomorphism numbers $(\hom(T,G))_T$.  Its zero set
$\{T:\varphi_T(G)=0\}=\{T: T\to G\}$ is the qualitative version of that vector.
Unlike the source-side quantity $|G,H|$, which differs for $2K_4$ and $K_4$
(both against $K_3$), the profile depends only on the homomorphism class of
$G$.

\begin{proposition}[Monotonicity and order embedding]
\label{prop:profile_order}
\begin{enumerate}
\item[(a)] If $G\to H$, then $\varphi_T(G)\ge\varphi_T(H)$ for every $T$.
      In particular, $\varphi(G)=\varphi(H)$ whenever $G$ and $H$ are
      homomorphically equivalent.
\item[(b)] If $\mathfrak F$ contains a graph homomorphically equivalent to
      $G$, then $\varphi(G)\ge\varphi(H)$ implies $G\to H$.
\end{enumerate}
Consequently, if $\mathfrak F$ contains all connected cores, then $G\mapsto\varphi(G)$
is an order-reversing embedding of the homomorphism order into
$([0,1]^{\mathfrak F},\le)$.  In that case, $E_w$ and $E_\infty$ are metrics on
homomorphism classes, and $\|\varphi(G)\|_1$ (for finite $\mathfrak F$) is a
strictly decreasing height function on the cores in $\mathfrak F$.
\end{proposition}

\begin{proof}
(a) If $T-S\to G\to H$, then $T-S\to H$.  This is the profile form of
\cite[Lemma~2.6]{efjt2015}.
(b) Let $T\in\mathfrak F$ with $T\leftrightarrow G$.  Then
$0=\varphi_T(G)\ge\varphi_T(H)$, so $G\to T\to H$.
Disconnected graphs are handled component by component.
\end{proof}

The embedding is not a lattice map.  The join and meet of the
homomorphism order are $G\sqcup H$ and $G\times H$.  These satisfy
$\varphi(G\sqcup H)\le\min(\varphi(G),\varphi(H))$ and
$\varphi(G\times H)\ge\max(\varphi(G),\varphi(H))$, and both inequalities can
be strict.  For example, take $T$ to be $K_3$ and the Gr\"otzsch graph $\mathrm{Gr}$
joined by a bridge.  Then $|T,K_3|=|T,\mathrm{Gr}|=1$ but
$|T,K_3\times\mathrm{Gr}|=2$.

\begin{proposition}[Clique and odd-cycle coordinates]
\label{prop:profile_closed}
Let $G$ have $\omega(G)=r\ge2$ and odd girth $g$ ($g=\infty$ if $G$ is
bipartite).  For every $n\ge2$ and $k\ge1$,
\[
  |K_n,G|=\binom n2-t\bigl(n,\min(n,r)\bigr),\qquad
  \varphi_{C_{2k+1}}(G)=
  \begin{cases}0,&2k+1\ge g,\\[2pt] \dfrac1{2k+1},&2k+1<g.\end{cases}
\]
Every bipartite test has $\varphi_T(G)=0$.  Hence $\omega(G)$ and the odd
girth of $G$ can be read off $\varphi(G)$:
$\omega(G)\ge n\iff\varphi_{K_n}(G)=0$, and
$g\le 2k+1\iff\varphi_{C_{2k+1}}(G)=0$.
\end{proposition}

\begin{proof}
The preserved edges of $K_n$ form a graph that maps to $G$, so this graph
is $K_{r+1}$-free.  By Tur\'an's theorem it has at most $t(n,r)$ edges.
Mapping $K_n$ in a balanced way onto an $r$-clique of $G$ attains this
bound.  For odd cycles, $C_{2k+1}\to G$ if and only if $G$ has an odd cycle
of length at most $2k+1$.  Otherwise, deleting one edge leaves a path,
which maps to an edge.
These are \cite[Props.~2.4, 2.5]{efjt2015} in the language of
separation.
\end{proof}

\begin{remark}
The argument improves Theorem~\ref{thm:lower_bound}: for all $G,H$,
\[
|G,H|\ \ge\ |E(G)|-t\bigl(|V(G)|,\omega(H)\bigr),
\]
because every subgraph of $G$ that maps to $H$ is
$K_{\omega(H)+1}$-free.  Since $\omega(H)\le\chi(H)$, this is never weaker.
It is attained by $G=K_n$ for \emph{every} $H$.
For example, $|K_5,C_5|=4$, while the $\chi$-bound gives only $10-t(5,3)=2$.
\end{remark}

\begin{proposition}[Range of the profile]
\label{prop:profile_range}
For every $G$ with an edge and every test $T$,
$\varphi_T(G)\le\varphi_T(K_{\omega(G)})\le 1/\omega(G)$.
Moreover $\varphi_{K_n}(G)\to1/\omega(G)$ as $n\to\infty$, so
$\sup_T\varphi_T(G)=1/\omega(G)$.  In particular,
$E_\infty(G,H)\ge\bigl|1/\omega(G)-1/\omega(H)\bigr|$.
\end{proposition}

\begin{proof}
The first inequality is Proposition~\ref{prop:profile_order}(a) applied to
$K_{\omega(G)}\to G$.  For the second, colour $V(T)$ uniformly at random
with $\omega(G)$ colours.  Each edge is properly coloured with probability
$1-1/\omega(G)$, so some colouring preserves at least that fraction of
$E(T)$.  The limit follows from Proposition~\ref{prop:profile_closed}.
\end{proof}

\begin{proposition}[Bounded tests cannot see $\chi$]
\label{prop:profile_blind}
Suppose every $T\in\mathfrak F$ has at most $m$ vertices, and $G$ has an
edge and odd girth larger than $m$.  Then $\varphi(G)=\varphi(K_2)$ on
$\mathfrak F$.
\end{proposition}

\begin{proof}
Let $f\colon V(T)\to V(G)$.  The image of the preserved edges is a subgraph
of $G$ on at most $m$ vertices.  It has no odd cycle, so it maps to an edge
of $G$.  Hence $\operatorname{mc}_G(T)=\operatorname{mc}_{K_2}(T)$.
\end{proof}

By Erd\H{o}s' theorem on graphs of large girth and chromatic number,
no finite family of tests detects $\chi$, $\chi_f$, or even
bipartiteness.  With summable weights, $E_w(G,K_2)\to0$ along any sequence
of graphs whose odd girth tends to infinity, whatever their chromatic
numbers.
The same phenomenon appears among small cores.  On the $1276$ cores
with at most $9$ vertices, every family of tests that separates all
homomorphism classes has at least $253$ members.  This is because $253$
pairs are separated by exactly one test, namely one of the two graphs
themselves.  A greedy search finds a separating family of size $311$.

\begin{proposition}[Relation to the separation of
F\"arnqvist--Jonsson--Thapper]
\label{prop:profile_fjt}
For edge-weighted graphs $(T,w)$, write
$p_G(T,w)=\operatorname{mc}_G(T,w)/\|w\|_1$; for unit weights,
$p_G(T)=1-\varphi_T(G)$.  The separation
of~\cite{fjt2009,efjt2015} is
\[
  s(G,H)=\inf_{(T,w)}\frac{p_G(T,w)}{p_H(T,w)}.
\]
\begin{enumerate}
\item[(a)] If $T$ is edge-transitive, then $\varphi_T(G)=1-s(G,T)$
      \cite[Thm.~2.1]{efjt2015}.
\item[(b)] The additive form of the FJT metric is the Hilbert projective
      metric on profiles:
      \[
        -\log s(G,H)-\log s(H,G)
        =\log\sup_{(T,w)}\frac{p_H(T,w)}{p_G(T,w)}
        +\log\sup_{(T,w)}\frac{p_G(T,w)}{p_H(T,w)}.
      \]
      For a finite unweighted family $\mathfrak F$, the same expression with
      the suprema over $T\in\mathfrak F$ is a pseudometric.  It is a metric on
      the homomorphism classes of graphs whose cores lie in $\mathfrak F$.
\end{enumerate}
\end{proposition}

\begin{proof}
(a) is the edge-transitive case of \cite[Thm.~2.1]{efjt2015}.
(b) is the definition of $s$, rewritten.  Positivity on a finite family
follows as in Proposition~\ref{prop:profile_order}(b).
\end{proof}

Thus the sup-distance $E_\infty$, the Euclidean distance $E_w$ and the FJT metric
are three distances on one embedding $G\mapsto\varphi(G)$.  Only the FJT
metric carries an approximation-theoretic meaning~\cite[Lemma~1.6]{efjt2015}.

\begin{proposition}[Moving the target; link to concentration]
\label{prop:profile_lip}
Let $S\subseteq E(T)$ with $|S|=|T,G|$, and let $f\colon T-S\to G$ be a
homomorphism.  For every graph $H$,
\[
  |T,H|\ \le\ |T,G|+\mu(f)\,|G,H|.
\]
In particular, if $T\to G$, then $|T,H|\le\gamma(T,G)\,|G,H|$.  Also
$|T,G-e|\le|T,G|+\mu(f,e)$ for every $e\in E(G)$.
\end{proposition}

\begin{proof}
Let $S'\subseteq E(G)$ be optimal with $g\colon G-S'\to H$.  Delete from
$T-S$ the at most $\mu(f)|S'|$ edges that $f$ maps into $S'$.  The map
$g\circ f$ is then a homomorphism from what remains to $H$.
\end{proof}

So the concentration parameter $\gamma$ is the local Lipschitz constant of
the profile under changes of the target.  The jump can be large:
deleting a triangle edge from $K_3$ with a long pendant path changes
$\varphi_{K_3}$ from $0$ to $1/3$, however many edges the graph has.

\begin{remark}[Computation]
$\operatorname{mc}_G(T)=\lim_{\beta\to\infty}\beta^{-1}\log\hom(T,W_\beta)$,
where $W_\beta$ is the looped complete graph on $V(G)$ with weight $e^\beta$
on $E(G)$ and $1$ elsewhere.  Consequently, the tree-decomposition dynamic
program for $\hom(T,\cdot)$, run over the $(\max,+)$ semiring, computes
$\varphi_T(G)$ in time $O(|V(T)|\,|V(G)|^{\operatorname{tw}(T)+1})$.
\end{remark}

% --- Suggested figure (data and plot by agents Q/R, discussion2/R_axes.png) ---
% \begin{figure}[ht]\centering
% \includegraphics[width=.9\linewidth]{figures/profile_pca.png}
% \caption{Principal components of $\varphi$ for the 85 cores ($\le 8$ vertices,
%  tests = the same cores). All Hasse covers of $\to$ point downward along
%  PC1 $\approx\|\varphi\|_1$ (Proposition~\ref{prop:profile_order}); PC2 is the
%  clique number (Proposition~\ref{prop:profile_closed}).}
% \label{fig:profile}
% \end{figure}

\subsection*{Open questions on the profile}
\begin{enumerate}
\item \textbf{Size of a separating family.}
  Let $\sigma(n)$ be the least size of a family $\mathfrak F$ such that $\varphi$
  restricted to $\mathfrak F$ separates all cores on at most $n$ vertices.
  We know $253\le \sigma(9)\le311$.  Is $\sigma(n)$ comparable to the number of
  cores on $n$ vertices, or is it much smaller?  Which cores are
  \emph{forced}, i.e.\ the only test separating some pair?
\item \textbf{The closure of the profile space.}
  The image of $\varphi$ is not closed: $\varphi(K_n)\to0$, and the
  Mycielski chain converges to a non-realised point.  Is the closure of
  $\{\varphi(G)\}$ in the product topology exactly the set of profiles of
  countable graphs?  Is the set of homomorphism classes totally bounded
  under $E_\infty$?
\item \textbf{Unweighted versus weighted separation.}
  Is
  $\sup_{X\to G}|X,T|/|E(X)|=1-s(T,G)$ for all $G,T$?  Here the supremum is
  over simple graphs $X$, and the right-hand side is the weighted FJT value.
  Equality holds when $G$ is edge-transitive (take $X=G$).  A positive
  answer would make both the FJT metric and the source-side profile
  definable without weights.
\end{enumerate}

% ---------------------------------------------------------------------------
% New bibitems (paste into thebibliography):
%
% \bibitem{fjt2009}
% T.~F\"arnqvist, P.~Jonsson and J.~Thapper,
% ``Approximability distance in the space of $H$-colourability problems,''
% in \textit{CSR 2009}, LNCS vol.~5675, pp.~92--104, 2009.
% doi:10.1007/978-3-642-03351-3\_11, arXiv:0802.0423.
%
% \bibitem{efjt2015}
% R.~Engstr\"om, T.~F\"arnqvist, P.~Jonsson and J.~Thapper,
% ``An approximability-related parameter on graphs -- properties and
% applications,''
% \textit{Discrete Mathematics \& Theoretical Computer Science},
% vol.~17, no.~1, pp.~33--66, 2015.
%
% \bibitem{lrs2006}
% M.~Langberg, Y.~Rabani and C.~Swamy,
% ``Approximation algorithms for graph homomorphism problems,''
% in \textit{APPROX 2006}, LNCS vol.~4110, pp.~176--187, 2006.
% doi:10.1007/11830924\_18.
%
% \bibitem{samal2015}
% R.~\v{S}\'amal,
% ``Cubical coloring -- fractional covering by cuts and semidefinite
% programming,''
% \textit{Discrete Mathematics \& Theoretical Computer Science},
% vol.~17, no.~2, 2015. arXiv:0911.2589.
%
% \bibitem{dgr2018}
% H.~Dell, M.~Grohe and G.~Rattan,
% ``Lov\'asz meets Weisfeiler and Leman,''
% in \textit{ICALP 2018}, LIPIcs vol.~107, pp.~40:1--40:14, 2018.
% doi:10.4230/LIPIcs.ICALP.2018.40.
%
% \bibitem{lovasz1967}   % [check]
% L.~Lov\'asz, ``Operations with structures,''
% \textit{Acta Math. Acad. Sci. Hungar.}, vol.~18, pp.~321--328, 1967.
% ---------------------------------------------------------------------------
